\begin{proof}[Proof of Lemma~\ref{lem:distinct}.]
    The centralizer $Z_\lambda$ of a permutation in $\C_\lambda$ is isomorphic, in this case, to the direct product $Z_\mu \times Z_\nu$. 
    By Definition~\ref{def:higher}, 
    $\omega^\lambda := \omega^{\mu} \otimes \omega^\nu$ and  
    \begin{equation}\label{eq:distinct}
        \psi^\lambda 
        = \omega^\lambda\uparrow_{Z_\lambda}^{S_n}
        = \left(\omega^{\mu}\uparrow_{Z_{\mu}}^{S_{|\mu|}} \otimes \, 
		\omega^{\nu}\uparrow_{Z_{\nu}}^{S_{|\nu|}}\right)
		\uparrow_{S_{|\mu|}\times S_{|\nu|}}^{S_n}.		    
    \end{equation}
    By the Littlewood-Richardson rule~\cite[Theorem A1.3.3]{EC2}, the outer product  of two irreducible characters 
    $\left( \chi^\alpha\otimes\chi^\beta\right)\uparrow_{S_m\times S_{n-m}}^{S_n}$ 
    contains irreducible representations indexed by hooks if and only if both $\alpha$ and $\beta$ are hooks; in the latter case,
    \[
        \langle \left( \chi^{(m-i,1^i)} \otimes \chi^{(n-m-j,1^j)}\right)\uparrow_{S_m \times S_{n-m}}^{S_n}, \chi^{(n-k,1^k)}\rangle = 
        \begin{cases}
            1, & \ k\in \{i+j, i+j+1\};\\
            0, & \ \rm{otherwise.}
        \end{cases}
    \]
    Therefore
    \[
        M_\lambda(x) = (1+x) M_\mu(x) M_\nu(x),
    \]
    as claimed.
\end{proof}	

\begin{corollary}
    If $\lambda$ is a partition with at least two different parts, namely not of the form $(r^s)$ for any $r$, 
    then $M_\lambda(x)$ is divisible by $1+x$ and the quotient has non-negative coefficients.
\end{corollary}

\begin{remark}\label{rem:distinct} 
    In this case, the existence of a cyclic descent extension
    may be proved directly as follows.
    By Equation~\eqref{eq:distinct}, $\psi^\lambda$ is a sum of characters indexed by disconnected shapes. 
    Thus, by 
    Corollary~\ref{cor:Schur-gf}, the distribution of the descent set over $\C_\lambda$ is equal to a sum of distributions over the sets of SYT of various disconnected shapes.
    By Theorem~\ref{conj1}, each of these sets carries a cyclic descent extension, hence so does~$\C_\lambda$.
\end{remark}


\section{Non-negativity: the single cycle case}
\label{sec:cycles}

Consider now the case of a conjugacy class with a single cycle.
By Corollary~\ref{cor:only if}, if $r$ is not square-free then $1+x$ divides $M_{(r)}(x)$. 
The main result of this section is the following. 

\begin{proposition}\label{prop:cycles_nonnegativity}
    If $r$ is not square-free then the coefficients of $M_{(r)}(x)/(1+x)$ are non-negative
\end{proposition}
		
It follows from Lemma~\ref{t:unimodal} below that, in order to prove Proposition~\ref{prop:cycles_nonnegativity}, it suffices to show the unimodality (to be defined) of $M_{(r)}(x)$. This is the content of the following statement.
		 
\begin{proposition}\label{prop:s=1_unimodal}
    For any positive integer $r$, the sequence $m_{0,(r)}, m_{1,(r)},\ldots, m_{r-1,(r)}$ is unimodal.
    The largest element is one of the middle ones, namely $m_{i,(r)}$ for $i = (r-1)/2$ if $r$ is odd, and either $i = (r-2)/2$ or $i = r/2$, or both, if $r$ is even.
\end{proposition}

In Subsection~\ref{sec:Witt-subsec} we use a variant of the Witt transform to produce explicit formulas for the coefficients $m_{j,(r)}$ (Lemma~\ref{t:m_formula}).
Then, in Subsection~\ref{sec:unimodality}, we prove their unimodality.


\subsection{A variant of the Witt transform}
\label{sec:Witt-subsec}
 
In this subsection we present a variant of the Witt transform, which will be used to prove non-negativity in Sections ~\ref{sec:unimodality} and~\ref{sec:Witt}.
    
Denote by $(i,j)$ the greatest common divisor of two integers $i,j$. Recall the arithmetical M\"obius function $\mu$. 
    
\begin{definition}\label{def:f}
    For a positive integer $r$ define 
    \[
        f_j(r) 
        := \frac{1}{r} \sum_{d|(r,j)} \mu(d)(-1)^{(d+1)j/d} \binom{r/d}{j/d} 
	   \qquad (0\le j\le r). 
    \]
\end{definition}
   
\begin{observation}\label{obs:f_values}
    By definition, 
    \[
        f_1(r) = f_{r-1}(r) = 1
        \qquad (r \ge 1).
    \]
    Also, the fundamental property
    \[
        \sum_{d|r} \mu(d) = 
        \begin{cases}
            1, &\text{if } r=1; \\
            0, &\text{if } r>1,
        \end{cases}
    \]
    and some case analysis ($r$ odd, or $r \equiv 0 \pmod 4$, or $r \equiv 2 \pmod 4$) imply that
    \[
        f_0(r) = 
        \begin{cases} 
            1,&\text{ if }  r=1;\\
            0,&\text{ if } r>1\\
        \end{cases}
        \qquad \text{and} \qquad
        f_r(r) = 
        \begin{cases} 
            1,&\text{ if }  r=1,2;\\
            0,&\text{ if } r>2.
        \end{cases}
    \]
\end{observation}
      
For the higher Lie character $\psi^{(r)}$, 
simplify slightly the notations in Equation~\eqref{e_k_interpretation1}:
\[
    m_{j,(r)}
    := \langle \psi^{(r)},\chi^{(r-j,1^j)} \rangle 
    \quad (0 \le j \le r-1), \quad
    e_{j,(r)} 
    :=\langle \psi^{(r)},\chi^{(1^j)\oplus(r-j)} \rangle
    \quad (0 \le j \le r).
\]

\begin{proposition}\label{prop:e_i_formuli1}
    For every $0\le j\le r$ 
    \[
        e_{j,(r)} = f_j(r).
    \]
    In particular, $f_j(r)$ is a non-negative integer.
\end{proposition}

\begin{remark}\label{rem:e=s}
    Proposition~\ref{prop:e_i_formuli1} will not be proved here, since it is the special case $s=1$ of Theorem~\ref{prop:e_i_formuli} below. 
    It also follows from a well known result of Kra\'skiewicz and Weyman~\cite{KraskiewiczWeyman} (Lemma~\ref{lem:Kra-Wey} below).
    A symmetric functions proof which applies~\cite[Lemma 2.7]{Sundaram_Adv} was presented by Sheila Sundaram~\cite{Sundaram_personal}. 
    Another proof follows from~\cite[Theorem 3.1]{ET}. 
    See Subsection~\ref{remarks_on_s=1} below for a discussion.
\end{remark}
  
\begin{definition}\label{def:F_polynomial}
    For a fixed positive integer $r$, collect the multiplicities $f_j(r)$ into a polynomial
    \[
	    F_r(x) := f_0(r) + f_1(r)x + f_2(r)x^2+ \ldots + f_r(r)x^r.
    \]
\end{definition}

Equation~\eqref{eq:alternating} and Proposition~\ref{prop:e_i_formuli1} imply the following. 

\begin{corollary}\label{t:F_and_M}
    \[
        F_r(x) = (1+x) M_{(r)}(x).
    \]
\end{corollary}

\begin{observation}\label{t:F}
    \[
        F_r(x) = \frac{1}{r} \sum_{d|r} \mu(d) (1-(-x)^d)^{r/d}.
    \]
\end{observation}

\begin{proof}
    Use Definition~\ref{def:f}, and write $j = kd$ if $d | (r,j)$. Then
    \begin{align*}
        F_r(x) 
        &= \sum_{j=0}^{r} x^j \sum_{d|(r,j)} \frac{\mu(d) (-1)^{(d+1)j/d}}{r} \binom{r/d}{j/d} 
        = \sum_{d|r} \frac{\mu(d)}{r} \sum_{k=0}^{r/d} (-1)^{(d+1)k} \binom{r/d}{k}x^{k d} \\
        &= \sum_{d|r} \frac{\mu(d)}{r} (1-(-x)^d)^{r/d}.
    \end{align*}

\end{proof}
	
\begin{remark}
    Recall from \cite{Moree} that the $r$-th Witt transform of a polynomial $p(x)$ is defined by 
    \[
        \mathcal{W}_p^{(r)}(x)=\frac{1}{r}\sum_{d|r}\mu(d)p(x^d)^{r/d}.
    \] 
    In our case, put $p(x)=1-x$ to get $F_r(x) = \mathcal{W}_p^{(r)}(-x)$. 
    The proof of Theorem~4 and Lemma~1 in~\cite{Moree} could have been used to prove that the coefficients of $F_r(x)$ are non-negative integers. This non-obvious property of the numbers $f_j(r)$ also follows, of course, from their interpretation in Proposition~\ref{prop:e_i_formuli1} as inner products of two characters.
    What we really need, in Proposition~\ref{prop:cycles_nonnegativity}, is the nonnegativity of the coefficients of~$F_r(x)/(1+x)^2$.
\end{remark}

We now produce an explicit formula for each coefficient $m_{j,(r)}$. For a combinatorial interpretation of these numbers, see Lemma~\ref{lem:Kra-Wey} below.

\begin{lemma}\label{t:m_formula}
    For a positive integer $r$,
    \[
		m_{j,(r)}
		= \frac{1}{r} \sum_{d|r} \mu(d) \binom{r/d - 1}{\lfloor j/d \rfloor}  (-1)^{j+\lfloor j/d \rfloor} 
		\qquad (0 \le j \le r-1).
    \]
\end{lemma}
\begin{proof}
	By Corollary~\ref{t:F_and_M},
	\[
	    (1+x) \sum_{j=0}^{r-1} m_{j,(r)}x^j=F_r(x). 
	\]
	
	Using Definition~\ref{def:f} and Observation~\ref{t:F},
	we can write
	\[
	    \sum_{j=0}^{r-1} m_{j,(r)} x^j 
	    = \frac{1}{r(1+x)} \sum_{d|r} \mu(d) (1-(-x)^d)^{r/d}.
    \]
	Using
	\[
		\frac{(1-(-x)^d)^{r/d}}{1+x} 
		= (1-(-x)^d)^{r/d-1} (1 - x + x^2 - \ldots +(-x)^{d-1})
	\]
	and comparing coefficients of $x^j$, where $j = dk + \ell$ with $0 \le \ell \le d-1$, we get
	\[
		r m_{j,(r)}
		= \sum_{d|r} \mu(d) \binom{r/d - 1}{k}  (-1)^{(d+1)k + \ell}
		= \sum_{d|r} \mu(d) \binom{r/d - 1}{\lfloor j/d \rfloor}  (-1)^{j+\lfloor j/d \rfloor}.
        \qedhere
	\]
\end{proof}


\subsection{Unimodality}
\label{sec:unimodality}
	
A sequence $(a_0,\ldots,a_n)$ of real numbers is called {\em unimodal} if there exists an index $0 \le i_0 \le n$ such that the sequence is weakly increasing up to position $i_0$ and weakly decreasing afterwards:
$a_0 \le a_1 \le \ldots \le a_{i_0} \ge \ldots a_{n-1} \ge a_n$.
	
\begin{lemma}\label{t:unimodal}
    Let $a(x) = a_0 + a_1 x + \ldots + a_n x^n$ be a polynomial with real, nonnegative and unimodal coefficients.
    Assume that $1+x$ divides $a(x)$, and let $b(x) := a(x)/(1+x)$.
    Then the coefficients of $b(x)$ are nonnegative.
\end{lemma}
	
\begin{proof}
    Let $b(x) = b_0 + \ldots + b_{n-1} x^{n-1}$.
    Then $a_0 = b_0$, $a_n = b_{n-1}$, and
    \begin{equation}\label{eq:unimodal1}
        a_i = b_{i-1} + b_i 
        \qquad (1 \le i \le n-1)
    \end{equation}
    Of course, divisibility of $a(x)$ by $1+x$ implies that
    \[
        \sum_{i=0}^{n} (-1)^i a_i = a(-1) = 0.
    \]
    Inverting~\eqref{eq:unimodal1} we get
    \begin{equation}\label{eq:unimodal2}
        b_i = \sum_{j=0}^i (-1)^{i-j} a_j
        \qquad (0 \le i \le n-1)
    \end{equation}
    and, similarly,
    \begin{equation}\label{eq:unimodal3}
        b_i = \sum_{j=i+1}^n (-1)^{j-i-1} a_j
        \qquad (0 \le i \le n-1).
    \end{equation}
    By assumption, the sequence $(a_0, \ldots, a_n)$ is nonnegative and unimodal, namely: 
    there exists an index $0 \le i_0 \le n$ such that
    \[
        0 \le a_0 \le \ldots \le a_{i_0} \ge \ldots \ge a_n \ge 0.
    \]
    It follows from~\eqref{eq:unimodal2} that, for odd indices $0 \le 2i+1 \le i_0$,
    \[
        b_{2i+1} = (a_{2i+1} - a_{2i}) + \ldots + (a_1 - a_0) \ge 0
    \]
    and, for even indices $0 \le 2i \le i_0$,
    \[
        b_{2i} = (a_{2i} - a_{2i-1}) + \ldots + (a_2 - a_1) + a_0 \ge 0.
    \]
    By~\eqref{eq:unimodal3}, a similar argument holds for 
    indices greater or equal to $i_0$, and the proof is complete.
\end{proof}
	
Lemma~\ref{t:unimodal} shows that non-negativity of a sequence can be proved by showing unimodality of a related sequence. In particular, Proposition~\ref{prop:cycles_nonnegativity} would follow once we show the unimodality of the polynomial $M_{(n)}(x)$.
In order to do that, we need the following technical lemma.

\begin{lemma}\label{lem:binomial_approx}
    Assume that $r > 7$ and $1 < j < r/2$;
    if $j = (r-1)/2$ assume also that $r > 11$.
    Let $d > 1$ be a divisor of $r$, and denote
	\[
	   A_{r,j,d} 
	   := \frac{\displaystyle\binom{r/d - 1} {\lfloor j/d \rfloor}} {\displaystyle\binom{r-1}{j}}.
	\]
	Then
	\[
	\frac{(r-1)(r-j)}{r-2j} A_{r,j,d} \le 
	\begin{cases}
	        1, &\text{if } d>2; \\
	        3/2, &\text{if } d=2.
	\end{cases}
	\]
\end{lemma}
	
\begin{proof}
    Write
	\[
	   \binom{r-1}{j}
	   = \frac{(r-1)(r-2) \cdots (r-j)}{j!}
	   = \prod_{1 \le i \le j} \frac{r-i}{i}.
	\]
	Let $\ell := \lfloor j/d \rfloor$. Then $\ell d$ is the largest multiple of $d$ not exceeding $j$, hence
	\[
	   \binom{r/d - 1}{\lfloor j/d \rfloor} 
	   = \binom{r/d - 1}{\ell}
	   = \frac{(r-d)(r-2d)\cdots (r-\ell d)}{d \cdot 2d \cdots \ell d}
	   = \prod_{1 \le i \le j,\, d \mid i} \frac{r-i}{i},
	\]
	The quotient $A_{r,j,d}$ can therefore be written in the form
	\begin{equation*}
		A_{r,j,d} 
		= \prod_{1 \le i \le j,\, d \nmid i} \frac{i}{r-i}.
	\end{equation*}
	By assumption $j < r/2$, thus $i/(r-i) < 1$ for all $1 \le i \le j$. It follows that $A_{r,j,d}$ is a decreasing function of $j$, with $A_{r,1,d} = 1/(r-1)$.
    
    For $d > 2$ and $2 \le j \le (r-2)/2$,
    \[
        A_{r,j,d} \le A_{r,2,d} = \frac{2}{(r-1)(r-2)}
        \le \frac{r-2j}{(r-1)(r-j)},
    \]
    where the last inequality, equivalent to $2(r-j) \le (r-2j)(r-2)$, follows from $2 \le r-2j$ and $r-j \le r-2$. 
	
	Similarly, for $d=2$ and $3 \le j \le (r-2)/2$,
	\[
	   A_{r,j,2} \le A_{r,3,2} = \frac{3}{(r-1)(r-3)} 
	   \le \frac{3(r-2j)}{2(r-1)(r-j)},
	\]
	where the last inequality, equivalent to $2(r-j) \le (r-2j)(r-3)$, follows from $2 \le r-2j$ and $r-j \le r-3$.

    For $d = 2$ and $j = 2$,
    \[
        A_{r,2,2} = \frac{1}{r-1} \le \frac{3(r-4)}{2(r-1)(r-2)},
    \]
    where the inequality, equivalent to $2(r-2) \le 3(r-4)$, follows from $r \ge 8$.
        
	It remains to consider the case $j = (r-1)/2$, namely $r = 2j+1$, for $d \ge 2$.
	Note that in this case we assumed that $r > 11$, namely $j > 5$.
	
	Assume first that $d > 2$. Then 
	\[
        A_{r,j,d} \le A_{r,6,d} \le A_{r,4,d}
        = \prod_{1 \le i \le 4,\, d \nmid i} \frac{i}{r-i}
        = \frac{a_d}{\displaystyle\binom{r-1}{4}},
    \]
    where
    \[
        a_d =
        \begin{cases}
            1, &\text{if } d > 4; \\
            (r-d)/d, &\text{if } d = 3,4.
        \end{cases}
    \]
    Clearly	
	\[
	   1 < \frac{r-4}{4} < \frac{r-3}{3} 
	\]
	and therefore 
	\[
	   A_{r,j,d} \le A_{r,4,d} \le A_{r,4,3} 
	   = \frac{8}{(r-1)(r-2)(r-4)}
	   \le \frac{r-2j}{(r-1)(r-j)},
	\]
	where the last inequality, equivalent (since $r = 2j+1$) to $8(j+1) \le (2j-1)(2j-3)$ and to $4j^2 - 16j \ge 5$, holds for $j \ge 5$.
	
	Finally, assume that $r = 2j+1$ and $d = 2$.
	Then 
	\[
	   A_{r,j,2} \le A_{r,6,2} = A_{r,5,2} 
	   = \frac{15}{(r-1)(r-3)(r-5)}
	   \le \frac{3(r-2j)}{2(r-1)(r-j)},
	\]
	where the last inequality, equivalent (since $r = 2j+1$) to $5(j+1) \le 2(j-1)(j-2)$ and to $2j^2 - 11j \ge 1$, holds for $j \ge 6$.
	This completes the proof.
\end{proof}
	
\begin{proof}[Proof of Proposition~\ref{prop:s=1_unimodal}.]
    We need to show that $m_{0,(r)} \le m_{1,(r)} \le \ldots \le m_{\lfloor (r-1)/2 \rfloor}$ and     $m_{r-1,(r)} \le m_{r-2,(r)} \le \ldots \le m_{\lceil (r-1)/2 \rceil}$ for any positive integer $r$.
    
	For $1 \le r \le 7$, computing the polynomial $M_r(x) := F_r(x)/(1+x)$ explicitly, using Observation~\ref{t:F}, gives
	\begin{align*}
		M_1(x) &= 1; \qquad 
		M_2(x) = M_3(x) = x; \qquad
		M_4(x) = x+x^2; \\
		M_5(x) &= x+x^2+x^3; \qquad
		M_6(x) = x+2x^2+x^3+x^4; \\
		M_7(x) &= x+2x^2+3x^3+2x^4+x^5.
	\end{align*}
	The claim clearly holds in these cases.
	Assume from now on that $r > 7$.
	
    Informally, the explicit formula for $m_{j,(r)}$ in Lemma~\ref{t:m_formula} has a dominant term corresponding to $d=1$, i.e.,
    $rm_{j,(r)}$ is approximately equal to $\binom{r - 1}{j}$. 
    We will show that this approximation is good enough to make the sequence $m_{0,(r)}, \ldots, m_{r-1,(r)}$ unimodal, like the sequence of binomial coefficients.
    Note that, unlike the binomial coefficients, this sequence is not always palindromic; see Proposition~\ref{conj:r2} below.
    
	We first show that $m_{j-1,(r)} \le m_{j,(r)}$ for $1 \le j < r/2$. Recall that we assume $r > 7$.
	For $j=1$, Lemma~\ref{t:m_formula} shows that, 
	for $r > 1$, $m_{0,(r)} = 0 < 1 = m_{1,(r)}$. 
	
    Assume now that  $1 < j < r/2$.	
    Clearly, for these values of $j$ and any divisor $d$ of $r$,
	\[
	   \binom{r/d - 1}{\lfloor j/d \rfloor} 
	   \ge \binom{r/d - 1}{\lfloor (j-1)/d \rfloor}.
	\]
	We conclude, by Lemma~\ref{t:m_formula}, that
	\[
	    rm_{j,(r)} - rm_{j-1,(r)} 
	    \ge \left[ \binom{r-1}{j} - \binom{r-1}{j-1} \right] - 2\sum_{d|r,\,d>1} \binom{r/d - 1}{\lfloor j/d \rfloor}. 
    \]
	Since
	\[
	    \binom{r-1}{j} - \binom{r-1}{j-1}
	    = \binom{r-1}{j} \left(1 - \frac{j}{r-j}\right)
	    = \binom{r-1}{j} \frac{r-2j}{r-j},
	\]
    using the notation of Lemma~\ref{lem:binomial_approx} we get
	\[
	   \frac{rm_{j,(r)} - rm_{j-1,(r)}} {\binom{r-1}{j}\frac{r-2j}{r-j}}
	   \ge 1 - 2\frac{r-j}{r-2j} \sum_{d|r,\,d>1} A_{r,j,d}.
	\]
	Let $d(r)$ denote the number of divisors of $r$. 
	For odd $r>7$ (unless $j = (r-1)/2$ and $r \in \{9,11\}$), Lemma~\ref{lem:binomial_approx} implies that
	\begin{align*}
		\frac{rm_{j,(r)} - rm_{j-1,(r)}} {\binom{r-1}{j}\frac{r-2j}{r-j}} 
		\ge 1 - \sum_{d|r,\,d>2} \frac{2}{r-1} 
		= 1 - \frac{2d(r)-2}{r-1}.
	\end{align*}
	For even $r>7$, Lemma~\ref{lem:binomial_approx} implies that
	\begin{align*}
		\frac{rm_{j,(r)} - rm_{j-1,(r)}} {\binom{r-1}{j}\frac{r-2j}{r-j}} 
		\ge 1 - \frac{3}{r-1} - \sum_{d|r,\,d>2} \frac{2}{r-1}
		= 1 - \frac{2d(r)-1}{r-1}.
	\end{align*}
    We clearly have $2d(r) \le r$ for $r>7$, and therefore $m_{j-1,(r)} \le m_{j,(r)}$ in both cases.
	
	In the remaining cases, namely $j = (r-1)/2$ and $r \in \{9,11\}$, we can compute directly using 	Lemma~\ref{t:m_formula}.
	For $r=9$ and $j=4$ the divisors are $d=1,3,9$, but $\mu(9) = 0$. Thus
	\[
		9m_{4,(9)}-9m_{3,(9)} 
		= \left[ \binom{8}{4} + \binom{2}{1} \right]
		- \left[ \binom{8}{3} - \binom{2}{1} \right]
		= 72 - 54 > 0.
	\]
	For $r=11$ and $j=5$ the divisors are $d=1,11$. Thus
	\[
		11m_{5,(11)}-11m_{4,(11)} 
		= \left[ \binom{10}{5} + \binom{0}{0} \right]
		- \left[ \binom{10}{4} - \binom{0}{0} \right]
		= 253 - 209 > 0.
	\]
    
	So far we have proven that $m_{0,(r)} \le m_{1,(r)} \le \ldots \le m_{\lfloor (r-1)/2 \rfloor, (r)}$ for $r > 7$.
	
	The remaining inequalities, $m_{r-1,(r)} \le m_{r-2,(r)} \le \ldots \le m_{\lceil (r-1)/2 \rceil, (r)}$, can be written as $m_{r - 1 - (j-1),(r)} \le m_{r - 1 - j,(r)}$ for $1 \le j < r/2$.
	By Lemma~\ref{t:m_formula},
	\[
		m_{r-1-j,(r)}
		= \frac{1}{r} \sum_{d|r} \mu(d) \binom{r/d - 1}{\lfloor (r-1-j)/d \rfloor}  (-1)^{r-1-j+\lfloor (r-1-j)/d \rfloor} 
		\qquad (0 \le j \le r-1).
    \]
    For a divisor $d$ of $r$, if $j = kd + \ell$ with $0 \le \ell \le d-1$ 
    then $r-1-j = (r/d-k-1)d + (d-1-\ell)$ with $0 \le d-1-\ell \le d-1$, 
    so that
    \[
        \lfloor j/d \rfloor + \lfloor (r-1-j)/d \rfloor
        = k + (r/d-k-1)) = r/d-1.
    \]
    It follows that
	\begin{align*}
		m_{r-1-j,(r)}
		&= \frac{1}{r} \sum_{d|r} \mu(d) \binom{r/d - 1}{\lfloor j/d \rfloor}  (-1)^{r-1-j+r/d-1-\lfloor j/d \rfloor} \\
		&= \frac{1}{r} \sum_{d|r} \mu(d) \binom{r/d - 1}{\lfloor j/d \rfloor}  (-1)^{r+r/d-j-\lfloor j/d \rfloor} 
		\qquad (0 \le j \le r-1).
    \end{align*}
    This is exactly the formula for $m_{j,(r)}$ except for the signs of the summands, which differ (for each $d|r$) by the factor $(-1)^{r+r/d}$.
    These signs do not play any role in the proof above that $m_{j-1,(r)} \le m_{j,(r)}$ for $1 \le j < r/2$, which therefore also shows, mutatis mutandis, that $m_{r-1-(j-1),(r)} \le m_{r-1-j,(r)}$ for $1 \le j < r/2$ --- except possibly the explicit confirmation when $j = (r-1)/2$ and $r \in \{9,11\}$. 
    In these cases $r$ is odd, and therefore $(-1)^{r+r/d} = 1$ for any divisor $d$ of $r$.
    This implies that indeed 
	\[
		m_{4,(9)}-m_{5,(9)} 
		= m_{4,(9)}-m_{3,(9)}
		> 0
	\]
	and
	\[
		m_{5,(11)}-m_{6,(11)} 
		= m_{5,(11)}-m_{4,(11)} 
		> 0,
	\]
    completing the proof.
\end{proof}

\begin{remark}
	We conjecture that Proposition~\ref{prop:s=1_unimodal} remains true for an arbitrary partition $\mu \vdash n$, in particular for every partition $(r^s) \vdash n$ with any	$s\ge 1$; see Conjecture~\ref{conj:unimodal}.
\end{remark}

\begin{proof}[Proof of Proposition~\ref{prop:cycles_nonnegativity}.]
    Combine Proposition~\ref{prop:s=1_unimodal} with   Lemma~\ref{t:unimodal}.
\end{proof}

We conclude 

\begin{corollary}\label{prop:CDE_cycles}
    The conjugacy class of $n$-cycles $\C_{(n)}$ carries a cyclic descent extension if and only if $n$ is not square-free.
\end{corollary}
	 
\begin{proof} 
    Combine Lemma~\ref{lem:2} with Corollary~\ref{cor:only if}(2) and Proposition~\ref{prop:cycles_nonnegativity}. 
\end{proof}


\section{Non-negativity: the case of one cycle length}
\label{sec:Witt}

In this section we consider the case $\lambda = (r^s)$. 
We fix $r$, while $s$ and hence $n=rs$ vary. 
The arguments below also work for the trivial case $r=1$.
		
As in the previous section, instead of the hook multiplicities 
\[
    m_{i,(r^s)} 
    = \langle \psi^{(r^s)}, \chi^{(n-i,1^i)} \rangle
\]
we prefer to work with their consecutive sums,
\[
    e_{i,(r^s)} := m_{i,(r^s)} + m_{i-1,(r^s)}.
\]
		
Here is the structure of the current section.
In Subsection~\ref{sec:e} 
we obtain an explicit description of $e_{i,(r^s)}$ (Theorem~\ref{prop:e_i_formuli}). The proof involves a detailed computation of character values and inner products of characters.
In Subsection~\ref{sec:prod} we 
transform this description into
a product formula (Corollary~\ref{power_series_form})
for the formal power series 
